% GENERATED FILE. Edit canonical lessons, then run npm run generate:books.
% canonical-source-hash: fnv1a64:d6c2ddf29273f2de
% canonical-lessons: FA-C95-hule, FA-C95-fenjan, FA-C95-dig, FA-C95-mahitabe, FA-C95-ojagh

\chapter{Towel, Cup, Pot, Frying pan, Stove}
\label{ch:fa-towel-cup-pot-frying-pan-stove}
\hlchaptermodality{95}

\begin{chapteropening}
\textbf{By the end of this chapter:} \emph{I can say a towel, a cup, a pot, a frying pan and a stove.}

Complete the last lesson of chapter 95: I can say a towel, a cup, a pot, a frying pan and a stove.
\end{chapteropening}

\section[\texorpdfstring{\fa{حوله} (hule)}{حوله (hule)}]{\fa{حوله} (hule) — a towel}

\label{lesson:FA-C95-hule}



\begin{quote}
Before the new one: say the Persian for thread, then the Persian for soap.
\end{quote}

\subsection*{You'll want to know: \fa{حوله}}
\textbf{\fa{حوله}} — \emph{hule} — \textquotedblleft{}a towel\textquotedblright{}.

\begin{cousinweb}
\emph{hule} is the everyday word for a towel; it is sometimes linked with Turkish \emph{havlu}, \textquotedblleft{}towel\textquotedblright{}, but its origin is not settled.
\end{cousinweb}

\begin{grammarlens}[title={What you've built}]
One more word for this chapter.
\end{grammarlens}

\subsection*{Your turn}
\begin{itemize}
\raggedright
  \item \emph{Say it:} \emph{hule}
  \item \emph{Say it:} \emph{hule}, once more
  \item \emph{Say it:} say \emph{hule}
  \item \emph{Recall:} say the Persian for thread, then the Persian for soap, then say \emph{hule} again
\end{itemize}

\begin{tcolorbox}[breakable,colback=teal!4,colframe=teal!35!black,title={Before you move on}]
What does \fa{حوله} mean? (\textquotedblleft{}A towel\textquotedblright{}.) Say it once more.
\end{tcolorbox}
\section[\texorpdfstring{\fa{فنجان} (fenjân)}{فنجان (fenjân)}]{\fa{فنجان} (fenjân) — a cup}

\label{lesson:FA-C95-fenjan}



\begin{quote}
Before the new one: say the Persian for soap, then the Persian for a towel.
\end{quote}

\subsection*{You'll want to know: \fa{فنجان}}
\textbf{\fa{فنجان}} — \emph{fenjân} — \textquotedblleft{}a cup\textquotedblright{}.

\begin{cousinweb}
\emph{fenjân} is usually explained as the Arabic form of the Persian \emph{pengân}, \textquotedblleft{}cup, bowl\textquotedblright{} (EtymDB): Arabic, with no \emph{p}, took the word as \emph{finjān}, and Persian took that form back. EtymDB traces \emph{pengân} to Greek \emph{pinax}, \textquotedblleft{}board, plate\textquotedblright{}.
\end{cousinweb}

\begin{grammarlens}[title={What you've built}]
One more word for this chapter.
\end{grammarlens}

\subsection*{Your turn}
\begin{itemize}
\raggedright
  \item \emph{Say it:} \emph{fenjân}
  \item \emph{Say it:} \emph{fenjân}, once more
  \item \emph{Say it:} say \emph{fenjân}
  \item \emph{Recall:} say the Persian for soap, then the Persian for a towel, then say \emph{fenjân} again
\end{itemize}

\begin{tcolorbox}[breakable,colback=teal!4,colframe=teal!35!black,title={Before you move on}]
What does \fa{فنجان} mean? (\textquotedblleft{}A cup\textquotedblright{}.) Say it once more.
\end{tcolorbox}
\section[\texorpdfstring{\fa{دیگ} (dig)}{دیگ (dig)}]{\fa{دیگ} (dig) — a pot}

\label{lesson:FA-C95-dig}



\begin{quote}
Before the new one: say the Persian for a towel, then the Persian for a cup.
\end{quote}

\subsection*{You'll want to know: \fa{دیگ}}
\textbf{\fa{دیگ}} — \emph{dig} — \textquotedblleft{}a pot\textquotedblright{}.

\begin{cousinweb}
\emph{dig} is a native Persian word, inherited rather than borrowed; it is usually traced to Middle Persian \emph{dīg}, \textquotedblleft{}cooking pot\textquotedblright{}.
\end{cousinweb}

\begin{grammarlens}[title={What you've built}]
One more word for this chapter.
\end{grammarlens}

\subsection*{Your turn}
\begin{itemize}
\raggedright
  \item \emph{Say it:} \emph{dig}
  \item \emph{Say it:} \emph{dig}, once more
  \item \emph{Say it:} say \emph{dig}
  \item \emph{Recall:} say the Persian for a towel, then the Persian for a cup, then say \emph{dig} again
\end{itemize}

\begin{tcolorbox}[breakable,colback=teal!4,colframe=teal!35!black,title={Before you move on}]
What does \fa{دیگ} mean? (\textquotedblleft{}A pot\textquotedblright{}.) Say it once more.
\end{tcolorbox}
\section[\texorpdfstring{\fa{ماهیتابه} (mâhitâbe)}{ماهیتابه (mâhitâbe)}]{\fa{ماهیتابه} (mâhitâbe) — a frying pan}

\label{lesson:FA-C95-mahitabe}



\begin{quote}
Before the new one: say the Persian for a cup, then the Persian for a pot.
\end{quote}

\subsection*{You'll want to know: \fa{ماهیتابه}}
\textbf{\fa{ماهیتابه}} — \emph{mâhitâbe} — \textquotedblleft{}a frying pan\textquotedblright{}.

\begin{cousinweb}
\emph{mâhitâbe} is \emph{mâhi}, \textquotedblleft{}fish\textquotedblright{}, with \emph{tâbe}, \textquotedblleft{}frying pan\textquotedblright{} (EtymDB): the pan for frying fish. Both parts are native; \emph{mâhi} continues Middle Persian \emph{māhīg} (EtymDB).
\end{cousinweb}

\begin{grammarlens}[title={What you've built}]
One more word for this chapter.
\end{grammarlens}

\subsection*{Your turn}
\begin{itemize}
\raggedright
  \item \emph{Say it:} \emph{mâhitâbe}
  \item \emph{Say it:} \emph{mâhitâbe}, once more
  \item \emph{Say it:} say \emph{mâhitâbe}
  \item \emph{Recall:} say the Persian for a cup, then the Persian for a pot, then say \emph{mâhitâbe} again
\end{itemize}

\begin{tcolorbox}[breakable,colback=teal!4,colframe=teal!35!black,title={Before you move on}]
What does \fa{ماهیتابه} mean? (\textquotedblleft{}A frying pan\textquotedblright{}.) Say it once more.
\end{tcolorbox}
\section[\texorpdfstring{\fa{اجاق} (ojâgh)}{اجاق (ojâgh)}]{\fa{اجاق} (ojâgh) — a stove}

\label{lesson:FA-C95-ojagh}



\begin{quote}
Before the new one: say the Persian for a pot, then the Persian for a frying pan.
\end{quote}

\subsection*{You'll want to know: \fa{اجاق}}
\textbf{\fa{اجاق}} — \emph{ojâgh} — \textquotedblleft{}a stove\textquotedblright{}.

\begin{cousinweb}
\emph{ojâgh} is a loan, usually traced to Turkic \emph{ocaq}, \textquotedblleft{}hearth\textquotedblright{}.
\end{cousinweb}

\begin{grammarlens}[title={What you've built}]
One more word for this chapter.
\end{grammarlens}

\subsection*{Your turn}
\begin{itemize}
\raggedright
  \item \emph{Say it:} \emph{ojâgh}
  \item \emph{Say it:} \emph{ojâgh}, once more
  \item \emph{Say it:} say \emph{ojâgh}
  \item \emph{Recall:} say the Persian for a pot, then the Persian for a frying pan, then say \emph{ojâgh} again
\end{itemize}

\begin{tcolorbox}[breakable,colback=teal!4,colframe=teal!35!black,title={Before you move on}]
What does \fa{اجاق} mean? (\textquotedblleft{}A stove\textquotedblright{}.) Say it once more.
\end{tcolorbox}
